\section{A range-preservation obstruction for full Gaussian path norm}
Minimizing the full Gaussian ReLU path norm can produce arbitrarily
large predictions between bounded observations. The following
construction concerns a prescribed finite sample; it makes no claim
about the probability of that sample geometry under an IID design.

For a univariate shallow ReLU implementation define
\[
 P(f)=\inf\left\{\sum_j|v_j|\sqrt{a_j^2+b_j^2}:
          f(x)=\sum_jv_j(a_jx+b_j)_+\ \text{on }[-1,1]\right\}.
\]
Every such representation satisfies
\[
 P(f)\ge\int_{(-1,1)}\sqrt{1+t^2}\,|f''|(dt), \tag{RG1}
\]
because an atom with knot $t=-b/a$ contributes second-derivative
mass $v|a|$, and
$\sqrt{1+t^2}|v a|=|v|\sqrt{a^2+b^2}$.
Cancellations only decrease the variation on the right.

Set $r=\sqrt3/2$ and $t_0=1/\sqrt3$. For
$0<\delta<1-r$, prescribe the four data pairs
\[
 (-r-\delta,-1),\quad(-r,0),\quad(r,0),\quad(r+\delta,-1).
\]
The three-neuron function
\[
 F_\delta(x)=\delta^{-1}
 \{r-t_0-(-x-t_0)_+-(x-t_0)_+\}
\]
interpolates them and has path norm at most
\[
 P(F_\delta)\le (r+\sqrt3)/\delta. \tag{RG2}
\]
Its middle plateau has height $(r-t_0)/\delta$.

More strongly, \emph{every} minimum-path-norm interpolant has
supremum magnitude at least $c/\delta$ for an absolute $c>0$.
To prove this, choose one fixed sufficiently small $\ell>0$ such that
\[
 a=\sqrt{1+(r-\ell)^2},\qquad 2a>r+\sqrt3.
\]
This is possible because
$2\sqrt{1+r^2}=\sqrt7>3\sqrt3/2=r+\sqrt3$.
Let $\psi_L$ be the triangular hat of height $a$, supported on
$[-r-\delta,-r+\ell]$, peaking at $-r$; let $\psi_R$ be its reflection.
Their supports are disjoint and
$0\le\psi_L+\psi_R\le\sqrt{1+x^2}$.
If an interpolant $f$ has $\|f\|_\infty\le M$, integration by parts gives
\[
 \int(\psi_L+\psi_R)\,df''
 =a\left[-\frac2\delta+
   \frac{f(-r+\ell)+f(r-\ell)}{\ell}\right].
\]
Consequently (RG1) implies
\[
 P(f)\ge \frac{2a}{\delta}-\frac{2aM}{\ell}. \tag{RG3}
\]
A minimum-path-norm interpolant has norm at most (RG2), so (RG3) yields
\[
 M\ge
 \frac{\ell[\,2a-(r+\sqrt3)\,]}{2a}\,\delta^{-1}.
\]
A minimizer exists: normalized signed unit-ridge evaluation vectors
form a compact subset of $\mathbb R^4$, so bounded-mass measures
have compact prediction sets. An optimum can be reduced to at most
four atoms by a linear-dependence argument.
Thus the statement concerns actual minimum-norm interpolants, not
only a particular nonminimal construction.

The example has a macroscopic unsampled central interval and two
very close pairs near its ends. Under a well-spread large IID sample
that geometry is rare. Moreover a likelihood with appreciable
regularization may decline to fit the close-pair label changes.
A successful posterior-specific improvement must quantify one of
those facts; a deterministic range-preservation assertion alone
cannot provide it.
